\documentclass[11pt]{article}
\usepackage[margin=1in]{geometry}
\usepackage{enumitem}
% =============================================================================
% Preambles/header.tex — shared LaTeX/Beamer preamble for this project
%
% Use from a Beamer lecture by prepending:
%     \input{header}
% and compiling with TEXINPUTS=../Preambles:$TEXINPUTS (the /compile-latex
% skill does this automatically).
%
% PALETTE CONTRACT. Color names defined here MUST match the names in
% Quarto/theme-template.scss so Beamer and Quarto renderings use the same
% palette. The scripts/check-palette-sync.sh script enforces this.
%
% To customize: edit the HEX values below AND the matching $variable in
% Quarto/theme-template.scss, then run scripts/check-palette-sync.sh.
% =============================================================================

% -----------------------------------------------------------------------------
% Engine detection + encoding
%
% Our /compile-latex skill uses XeLaTeX, where inputenc is unnecessary and
% can produce encoding warnings. Only load inputenc under pdfTeX.
% -----------------------------------------------------------------------------
\usepackage{iftex}
\ifPDFTeX
  \usepackage[utf8]{inputenc}
\fi

\usepackage{amsmath, amssymb, amsthm}
\usepackage{booktabs}
\usepackage{graphicx}

% -----------------------------------------------------------------------------
% PALETTE — mirrors Quarto/theme-template.scss variable names.
% Load xcolor and define colors BEFORE anything that references them
% (notably hyperref, hypersetup, and the Beamer color assignments).
% -----------------------------------------------------------------------------
\usepackage{xcolor}

\definecolor{primary-blue}{HTML}{012169}      % headings, accents
\definecolor{primary-gold}{HTML}{B9975B}      % emphasis, borders
\definecolor{highlight-yellow}{HTML}{F2A900}  % markers, alerts
\definecolor{light-bg}{HTML}{E8EDF5}          % title / section backgrounds
\definecolor{jet}{HTML}{1A1A1A}               % body text

% Semantic colors (used in diagrams and annotations)
\definecolor{positive}{HTML}{15803D}          % good / observed / identified
\definecolor{negative}{HTML}{B91C1C}          % bad / problematic / bias
\definecolor{neutral}{HTML}{525252}           % reference / context

% Highlight accents (match .hi-* classes in the SCSS)
\definecolor{hi-slate}{HTML}{314F4F}
\definecolor{hi-green}{HTML}{2E8B57}
\definecolor{hi-red}{HTML}{C41E3A}

% Semantic aliases so skill templates and snippets can write
% \textcolor{accent}{...} without hard-coding "primary-blue".
\colorlet{accent}{primary-blue}
\colorlet{accent2}{primary-gold}

% -----------------------------------------------------------------------------
% Hyperref — loaded AFTER xcolor + palette so link colors resolve.
% -----------------------------------------------------------------------------
\usepackage{hyperref}
\hypersetup{colorlinks=true, linkcolor=primary-blue, urlcolor=primary-blue,
            citecolor=primary-blue, pdfborder={0 0 0}}

% -----------------------------------------------------------------------------
% Course code — set once per deck (after \input{header}, before \begin{document})
% via \coursecode{CS401}. Shown in the footer on every slide so a deck's course
% is identifiable without opening the title page. Empty by default.
% -----------------------------------------------------------------------------
\makeatletter
\newcommand{\coursecode}[1]{\gdef\@coursecodetext{#1}}
\gdef\@coursecodetext{}
\makeatother

% -----------------------------------------------------------------------------
% Beamer theme assignments (only apply under Beamer).
%
% We detect Beamer via \@ifclassloaded{beamer}, which is the documented,
% non-internal way. This must be wrapped in \makeatletter/\makeatother
% because \@ifclassloaded contains an @-catcode character.
% -----------------------------------------------------------------------------
\makeatletter
\@ifclassloaded{beamer}{%
  \usefonttheme{professionalfonts}
  \setbeamercolor{structure}{fg=primary-blue}
  \setbeamercolor{frametitle}{fg=primary-blue}
  \setbeamercolor{title}{fg=primary-blue}
  \setbeamercolor{subtitle}{fg=primary-gold}
  \setbeamercolor{author}{fg=primary-blue}
  \setbeamercolor{institute}{fg=neutral}
  \setbeamercolor{date}{fg=primary-gold}
  \setbeamercolor{itemize item}{fg=primary-blue}
  \setbeamercolor{itemize subitem}{fg=primary-gold}
  \setbeamercolor{alerted text}{fg=primary-gold}
  \setbeamercolor{block title}{fg=white, bg=primary-blue}
  \setbeamercolor{block body}{bg=light-bg}
  % Semantic block variants: block=definition/concept (blue), exampleblock=
  % worked example (green/positive), alertblock=key takeaway or caution (gold).
  % Keep to <=2 colored boxes per slide across ALL three variants (INV-7).
  \setbeamercolor{block title example}{fg=white, bg=positive}
  \setbeamercolor{block body example}{bg=light-bg}
  \setbeamercolor{block title alerted}{fg=white, bg=primary-gold}
  \setbeamercolor{block body alerted}{bg=light-bg}

  % Minimal footer: course code (left) + page X / N (right), no navigation clutter
  \setbeamertemplate{navigation symbols}{}
  \setbeamertemplate{footline}{%
    \color{neutral}\footnotesize\hspace{0.7em}\@coursecodetext\hfill
    \insertframenumber/\inserttotalframenumber\hspace{0.7em}\vspace{0.3em}}
}{}
\makeatother

% -----------------------------------------------------------------------------
% TikZ setup — libraries the snippet gallery uses
% -----------------------------------------------------------------------------
\usepackage{tikz}
\usetikzlibrary{arrows.meta, positioning, calc, decorations.pathreplacing,
                fit, shapes.geometric, backgrounds}

% Shared TikZ styles — snippets and hand-written diagrams can pick these up
% via `\begin{tikzpicture}[every node/.style={...}]` or by naming specific
% styles (e.g., `\node[dag-node] (X) at (0,0) {X};`).
\tikzset{
  % Node styles for causal DAGs and similar
  dag-node/.style = {
    draw=primary-blue, fill=light-bg, rounded corners=2pt,
    minimum width=1.2cm, minimum height=0.9cm, align=center, font=\small
  },
  decision-node/.style = {
    draw=primary-gold, fill=white, diamond, aspect=1.8,
    minimum width=1.8cm, minimum height=1.2cm, align=center, font=\small,
    inner sep=1pt
  },
  % Edge styles — observed vs counterfactual
  observed-edge/.style = {-{Latex[length=2.5mm]}, thick, draw=jet},
  counterfactual-edge/.style = {-{Latex[length=2.5mm]}, thick, draw=neutral, dashed},
  confound-edge/.style = {-{Latex[length=2.5mm]}, thick, draw=negative, dashed},
  % Data-point styles for plots
  observed-dot/.style = {fill=primary-blue, draw=primary-blue, circle, inner sep=1.2pt},
  counterfactual-dot/.style = {fill=white, draw=neutral, circle, inner sep=1.2pt}
}

% -----------------------------------------------------------------------------
% Convenience macros
% -----------------------------------------------------------------------------
\newcommand{\muted}[1]{\textcolor{neutral}{#1}}
\newcommand{\key}[1]{\textcolor{primary-gold}{\textbf{#1}}}
\newcommand{\good}[1]{\textcolor{positive}{#1}}
\newcommand{\bad}[1]{\textcolor{negative}{#1}}

% Standout frame macro: full-bleed dark background for section transitions.
% Beamer-only (uses \begin{frame}/\setbeamercolor/\usebeamerfont) -- do not
% call from an article-class file (e.g. Notes/<CODE>/*-notes.tex). \muted,
% \key, \good, \bad, and \coursecode above are plain xcolor/gdef and are
% safe to use from either class; verified when Notes/article-class support
% was added.
% Expand: \transitionslide{Your transition title}
\newcommand{\transitionslide}[1]{%
  \begingroup
  \setbeamercolor{background canvas}{bg=primary-blue}
  \begin{frame}[plain]
    \centering
    \vfill
    {\usebeamerfont{title}\color{white}\Large #1\par}
    \vfill
  \end{frame}
  \endgroup
}

% Section divider frame (Beamer-only), an alternative to \transitionslide with a
% darker two-line style: near-black (jet) background, a small white label line
% above a larger gold title, and a thin gold rule along the bottom edge.
% Expand: \sectiondivider{Section 2}{Pointers}
\newcommand{\sectiondivider}[2]{%
  \begingroup
  \setbeamercolor{background canvas}{bg=jet}
  \begin{frame}[plain]
    \vspace*{3.0cm}
    {\color{white}\ttfamily\large #1\par}
    \vspace{0.5cm}
    {\color{primary-gold}\LARGE #2\par}
    \begin{tikzpicture}[remember picture, overlay]
      \draw[primary-gold, line width=2.5pt]
        ($(current page.south west)+(0,0.1)$) -- ($(current page.south east)+(0,0.1)$);
    \end{tikzpicture}
  \end{frame}
  \endgroup
}

% -----------------------------------------------------------------------------
% Channel watermark — "Concepts That Click" (YouTube).
%
% Article-class files (CompetitiveExam Q/A sets, Notes, Assignments) get a
% light diagonal draft watermark on every page plus a centered footer naming
% the channel. Beamer decks get a subtle TikZ background watermark instead
% (the footline already carries the course code). Centralized here so every
% MCQ PDF is branded without per-file edits.
% -----------------------------------------------------------------------------
\makeatletter
\@ifclassloaded{article}{%
  \usepackage{draftwatermark}
  \SetWatermarkText{Concepts That Click}
  \SetWatermarkScale{1}
  \SetWatermarkFontSize{2cm}
  \SetWatermarkLightness{0.86}
  \SetWatermarkAngle{45}
  \usepackage{fancyhdr}
  \pagestyle{fancy}
  \fancyhf{}
  \fancyfoot[C]{\footnotesize\textcolor{neutral}{Concepts That Click~|~YouTube: \href{https://www.youtube.com/@concepts.thatclick}{@concepts.thatclick}}}
  \renewcommand{\headrulewidth}{0pt}
  \renewcommand{\footrulewidth}{0pt}
  \fancypagestyle{plain}{%
    \fancyhf{}%
    \fancyfoot[C]{\footnotesize\textcolor{neutral}{Concepts That Click~|~YouTube: \href{https://www.youtube.com/@concepts.thatclick}{@concepts.thatclick}}}%
    \renewcommand{\headrulewidth}{0pt}%
    \renewcommand{\footrulewidth}{0pt}%
  }
}{}
\@ifclassloaded{beamer}{%
  \setbeamertemplate{background}{%
    \begin{tikzpicture}[remember picture, overlay]
      \node[rotate=45, opacity=0.055, align=center]
        at (current page.center)
        {\Huge\textcolor{neutral}{Concepts That Click}};
    \end{tikzpicture}%
  }
}{}
\makeatother 
\coursecode{JKSSB}
\title{Lesson 61 Practice: Current Office, Email and Statement Evaluation}
\author{JKSSB Computer Awareness}
\date{\today}
\begin{document}
\maketitle
\noindent\textbf{Provenance:} All 20 questions are original JKSSB-pattern
questions (\textit{[Original]}); concepts drawn from PYQ-12/13/14/15/16/17;
none reproduces an official paper item.

\begin{enumerate}[label=\textbf{Q\arabic*.}, leftmargin=*]
\item \textit{[Original]} In MS Word, content controls (drop-down list, date
  picker, check box) are inserted from which tab?
  \begin{enumerate}[label=\Alph*.]
    \item Developer tab
    \item Review tab
    \item Mailings tab
    \item References tab
  \end{enumerate}
\item \textit{[Original]} Which of the following statements about MS Word is
  \textbf{NOT} correct?
  \begin{enumerate}[label=\Alph*.]
    \item A content control can enforce limited data validation, such as a
      drop-down list or a date picker.
    \item A hanging indent shifts every line of the paragraph except the first
      line.
    \item A reviewer's comments on a sentence appear directly in the Navigation
      Pane.
    \item Accepting a tracked change keeps the new text and removes the
      revision metadata attached to that change.
  \end{enumerate}
\item \textit{[Original]} Which pairing of indent type and effect is correctly
  matched?
  \begin{enumerate}[label=\Alph*.]
    \item First-line indent --- indents every line except the first
    \item Hanging indent --- indents only the first line of the paragraph
    \item First-line indent --- indents only the first line of the paragraph
    \item Hanging indent --- centres the first line and left-aligns the rest
  \end{enumerate}
\item \textit{[Original]} The Word Navigation Pane offers which three views?
  \begin{enumerate}[label=\Alph*.]
    \item Headings, Pages, Results
    \item Comments, Footnotes, Endnotes
    \item Draft, Outline, Print Layout
    \item Insert, Design, References
  \end{enumerate}
\item \textit{[Original]} When Track Changes is switched on, each recorded
  change stores:
  \begin{enumerate}[label=\Alph*.]
    \item only the page number of the edit.
    \item the author's name and the date and time of the edit.
    \item only the final text with no author information.
    \item the printer name and the paper size.
  \end{enumerate}
\item \textit{[Original]} A report needs no header on the title page but a
  running header on every later page. Which break makes this possible?
  \begin{enumerate}[label=\Alph*.]
    \item Page Break
    \item Column Break
    \item Section Break
    \item Line Break
  \end{enumerate}
\item \textit{[Original]} Given
  \texttt{=IF(A1>50,"High",IF(A1>25,"Medium","Low"))} with A1 containing 30,
  the formula returns:
  \begin{enumerate}[label=\Alph*.]
    \item High
    \item Medium
    \item Low
    \item 0
  \end{enumerate}
\item \textit{[Original]} With the same formula
  \texttt{=IF(A1>50,"High",IF(A1>25,"Medium","Low"))}, if A1 is blank the
  result is:
  \begin{enumerate}[label=\Alph*.]
    \item High, because a blank cell is read as 100.
    \item Medium, because a blank cell skips the first test.
    \item Low, because a blank cell is treated as 0 in the comparison.
    \item \texttt{\#VALUE!}, because a blank cell always causes an error.
  \end{enumerate}
\item \textit{[Original]} Cell A1 contains the text ``abc''. Which statement
  is correct?
  \begin{enumerate}[label=\Alph*.]
    \item \texttt{=A1+5} returns 5.
    \item \texttt{ISNUMBER(A1)} returns TRUE.
    \item \texttt{=A1+5} returns \texttt{\#VALUE!} because text is not a
      number.
    \item Text is treated as 0 in every arithmetic operation.
  \end{enumerate}
\item \textit{[Original]} Which is the correct general form of COUNTIFS for
  counting rows that satisfy all the given conditions?
  \begin{enumerate}[label=\Alph*.]
    \item \texttt{COUNTIFS(criteria1, range1, criteria2, range2)}
    \item \texttt{COUNTIFS(range1, criteria1, range2, criteria2)}
    \item \texttt{COUNTIFS(range1 + range2, criteria1 AND criteria2)}
    \item \texttt{COUNTIFS(range1; criteria1; TRUE)}
  \end{enumerate}
\item \textit{[Original]} Which formula counts the rows in which A1:A10 exceeds
  50 \textbf{and} B1:B10 contains ``Yes''?
  \begin{enumerate}[label=\Alph*.]
    \item \texttt{=SUM(A1:A10>50, B1:B10="Yes")}
    \item \texttt{=SUMPRODUCT((A1:A10>50)*(B1:B10="Yes"))}
    \item \texttt{=AVERAGE(A1:A10, B1:B10)}
    \item \texttt{=COUNT(A1:A10, B1:B10)}
  \end{enumerate}
\item \textit{[Original]} Evaluate the statements:
  \begin{enumerate}[label=\Roman*.]
    \item COUNTIFS counts only the rows satisfying all the given
      range-and-criteria pairs.
    \item SUMPRODUCT can only add numbers and cannot count rows meeting
      multiple conditions.
    \item In a numeric comparison, a blank cell is treated as zero.
  \end{enumerate}
  \begin{enumerate}[label=\Alph*.]
    \item I and III only
    \item II only
    \item I, II, and III
    \item I and II only
  \end{enumerate}
\item \textit{[Original]} An email rule (filter) that moves all mail from one
  sender into a dedicated folder consists of:
  \begin{enumerate}[label=\Alph*.]
    \item a condition plus an action, applied automatically to mail.
    \item a digital signature plus encryption, applied manually.
    \item a Bcc field plus a Cc field, applied once.
    \item a Section Break plus a Page Break, applied to pages.
  \end{enumerate}
\item \textit{[Original]} Which statement correctly distinguishes a digital
  signature from encryption?
  \begin{enumerate}[label=\Alph*.]
    \item A digital signature encrypts the message body so only the recipient
      can read it.
    \item Encryption proves the sender's identity; a digital signature hides
      the content.
    \item A digital signature provides authenticity and integrity, while
      encryption provides confidentiality.
    \item A digital signature and encryption are different names for the same
      operation.
  \end{enumerate}
\item \textit{[Original]} Expand POP3.
  \begin{enumerate}[label=\Alph*.]
    \item Point-to-Point Protocol version 3
    \item Post Office Protocol version 3
    \item Packet Ordering Protocol version 3
    \item Private Office Protocol version 3
  \end{enumerate}
\item \textit{[Original]} Which statement about mail protocols is correct?
  \begin{enumerate}[label=\Alph*.]
    \item POP3 sends mail; SMTP downloads mail to the local device.
    \item IMAP downloads mail and always deletes it from the server; POP3
      keeps devices synchronised.
    \item SMTP sends mail; POP3 downloads mail to the local device; IMAP keeps
      mail on the server and synchronises devices.
    \item POP3 keeps mail synced on the server across all devices.
  \end{enumerate}
\item \textit{[Original]} A message is sent with addresses in the To, Cc, and
  Bcc fields. Who can see the Bcc addresses?
  \begin{enumerate}[label=\Alph*.]
    \item Every recipient, since Bcc is visible to all.
    \item Only the sender; the Bcc addresses are hidden from the other
      recipients.
    \item Only the Cc recipients; the sender cannot see them.
    \item Nobody, since Bcc addresses are deleted before sending.
  \end{enumerate}
\item \textit{[Original]} In PowerPoint, which content type CANNOT be converted
  to a SmartArt graphic in one step without re-entry?
  \begin{enumerate}[label=\Alph*.]
    \item A bulleted list of text
    \item A picture inserted on the slide
    \item Plain paragraph text with levels
    \item A numbered list of points
  \end{enumerate}
\item \textit{[Original]} Evaluate the statements:
  \begin{enumerate}[label=\Roman*.]
    \item A Section Break, not a Page Break, allows different headers and
      footers in one document.
    \item Addresses in the Bcc field are hidden from the sender.
    \item A digital signature encrypts the message body for confidentiality.
  \end{enumerate}
  \begin{enumerate}[label=\Alph*.]
    \item I only
    \item I and II only
    \item II and III only
    \item I, II, and III
  \end{enumerate}
\item \textit{[Original]} Which pairing is correctly matched?
  \begin{enumerate}[label=\Alph*.]
    \item Page Break --- starts a new section with its own headers and
      footers
    \item COUNTIFS --- handles one condition only; SUMPRODUCT --- only adds
      numbers
    \item First-line indent --- shifts only the first line; Hanging indent
      --- shifts every line except the first
    \item POP3 --- syncs mail on the server; IMAP --- downloads and deletes
      from the server
  \end{enumerate}
\end{enumerate}
\end{document}
